\documentclass[11pt]{article}
\usepackage{amsmath,amssymb,amsthm}

\newtheorem{theorem}{Theorem}
\newtheorem{corollary}[theorem]{Corollary}
\newtheorem{lemma}[theorem]{Lemma}

\newcommand{\Up}{U_+}
\newcommand{\mass}{\operatorname{B}}
\newcommand{\car}{\operatorname{C}}
\newcommand{\sh}{\operatorname{D}}

\title{Variable-depth inverse reduction and the first two decoder layers}
\author{}
\date{27 August 2026}

\begin{document}
\maketitle

For a positive odd integer $x$, put
\[
 \Up(x):=\frac{3x+1}{2^{v_2(3x+1)}}.
 \tag{1}
\]
For a positive exponent word $w=(a_1,\ldots,a_m)$, define
\[
 \mass(w):=\sum_{i=1}^m a_i,
 \qquad
 \car(w):=\sum_{i=1}^m3^{m-i}2^{a_1+\cdots+a_{i-1}},
 \tag{2}
\]
and its shifted carry by
\[
 \sh(w):=\car(w)+2^{\mass(w)}-3^m.
 \tag{3}
\]
Thus an exact path with word $w$ from $x$ to $y$ satisfies
\[
 2^{\mass(w)}(y+1)=3^m(x+1)+\sh(w).
 \tag{4}
\]
We use the record atoms
\[
 A:=1,\qquad B:=12,\qquad C:=1113,\qquad D:=1122.
 \tag{5}
\]

\section{A sharp prefix interval}

\begin{lemma}[shifted-carry interval]
Let $E=(e_1,\ldots,e_h)$ be a positive exponent word and write
\[
 \mass(E)=h+t,\qquad t\geq0.
 \tag{6}
\]
Then
\[
 0\leq\sh(E)\leq3^h(2^t-1).
 \tag{7}
\]
Equality on the right occurs if and only if $t=0$, in which case
$E=1^h$ and $\sh(E)=0$.

If $E$ is an exact path from a positive odd $z$ to a positive odd $r$,
put
\[
 \mathcal A_h(r):=\left(\frac23\right)^h(r+1).
 \tag{8}
\]
Then
\[
 \boxed{
  2^t\bigl(\mathcal A_h(r)-1\bigr)
  \leq z\leq2^t\mathcal A_h(r)-1.}
 \tag{9}
\]
The lower bound in (9) is attained if and only if $t=0$.
\end{lemma}

\begin{proof}
Expanding (3) one letter at a time gives
\[
 \sh(E)=\sum_{i=1}^h
 3^{h-i}2^{e_1+\cdots+e_{i-1}}(2^{e_i}-2)\geq0.
 \tag{10}
\]
Moreover,
\[
 \frac{\car(E)}{2^{\mass(E)}}
 =\sum_{i=1}^h\frac{3^{h-i}}{2^{e_i+\cdots+e_h}}
 \leq\sum_{j=0}^{h-1}\frac{3^j}{2^{j+1}}
 =\left(\frac32\right)^h-1.
 \tag{11}
\]
The inequality in (11) is an equality precisely when every $e_i=1$.
Substitution in (3) proves (7), including its equality assertion.

Equation (4) for $E$ is
\[
 z+1=2^t\mathcal A_h(r)-\frac{\sh(E)}{3^h}.
 \tag{12}
\]
The two bounds in (7) now give (9).  Equality at the lower edge forces
equality in (7), hence $t=0$; the converse follows from (10).
\end{proof}

\section{The exact variable-depth factorization}

We use the arbitrary-prefix same-depth inverse theorem: if
\[
 W=PU,\qquad P\in\{B,C,D\},\qquad U\in\{A,B,C,D\}^*,
 \tag{13}
\]
is exact from $q$ to $y$ and has raw length $L$, then its least positive
odd $L$-step predecessor is
\[
 q_*:=\begin{cases}
 q,&P=B\text{ or }D,\\[1mm]
 (q-1)/2,&P=C.
 \end{cases}
 \tag{14}
\]
It is unique; in the second case its word is $1121U$.

The same theorem gives a useful refinement.  If $F$ is any exact
length-$L$ word from $r$ to $y$, write
\[
 \delta:=\mass(F)-\mass(W),
 \qquad r=2^\delta q+c.
 \tag{15}
\]
Its distortion proof gives $\delta\geq-1$ and
$-1\leq c/2^\delta\leq5$.  At defect $-1$ one has exactly the second
case of (14).  At defect zero, parity gives $c\in\{0,2,4\}$.  In fact
$c=4$ is impossible: comparison of the two shifted equations gives
\[
 \sh(F)=\sh(W)-c3^L,
 \tag{16}
\]
where $0\leq\sh(F)$ and $\sh(W)<4\,3^L$.  Consequently every $F$ lies
in exactly one of the following channels:
\[
\begin{array}{c|c|c}
\text{defect}&\text{possible start}&\text{equality data}\\ \hline
-1&(q-1)/2&P=C,\quad F=1121U,\\
0&q\text{ or }q+2&r=q\Longleftrightarrow F=W,\\
\delta\geq1&r\geq2^\delta(q-1)&{}.
\end{array}
\tag{17}
\]

\begin{theorem}[variable-depth channel reduction]
Let $W=PU$ be as in (13), and let $V$ be any exact word of length
$L+h$ from a positive odd integer $z$ to $y$, where $h\geq0$.
Split it uniquely as
\[
 V=EF,
 \qquad |E|=h,
 \qquad |F|=L,
 \tag{18}
\]
and put
\[
 r:=\Up^h(z),
 \qquad t:=\mass(E)-h.
 \tag{19}
\]
Then $t\geq0$, the suffix $(F,r)$ belongs to exactly one of the three
channels (17), and
\[
 \boxed{
  2^t\bigl(\mathcal A_h(r)-1\bigr)
  \leq z\leq2^t\mathcal A_h(r)-1.}
 \tag{20}
\]
In particular,
\[
 \boxed{
 z\geq2^t\bigl(\mathcal A_h(q_*)-1\bigr).}
 \tag{21}
\]

Put $b_*:=v_3(q_*+1)$.  For every $0\leq h\leq b_*$, the unique least
positive odd $(L+h)$-step predecessor of $y$ is
\[
 \boxed{
 z_h+1=\left(\frac23\right)^h(q_*+1).}
 \tag{22}
\]
Its word is
\[
 \boxed{
 \begin{cases}
 1^hW,&P=B\text{ or }D,\\
 1^h1121U,&P=C.
 \end{cases}}
 \tag{23}
\]
\end{theorem}

\begin{proof}
The split (18) is made after the first $h$ exact steps of $V$; hence
$E$ is exact from $z$ to $r$ and $F$ is exact from $r$ to $y$.  The
channel statement follows from (17), while (20) is Lemma~1 applied to
$E$.  Since every channel has $r\geq q_*$, monotonicity in $r$ proves
(21).

If $h\leq b_*$, the right side of (22) is an even positive integer.
Starting from $z_h$, the first $h$ exponent-one steps are exact and end
at $q_*$.  Appending the unique minimum suffix in (14) gives (23).
This path has $t=0$ and $r=q_*$, so it attains (21).  Equality in (21)
forces equality both in $r\geq q_*$ and in Lemma~1; hence $r=q_*$,
$t=0$, $E=1^h$, and the suffix is the unique word in (14).  This proves
the asserted uniqueness.
\end{proof}

\begin{corollary}[necessary low-excess condition for a new descent]
In the notation of Theorem~2, suppose $\mathcal A_h(q_*)>1$.  If
$z<N$ for some positive threshold $N$, then necessarily
\[
 \boxed{
 2^t<\frac{N}{\mathcal A_h(q_*)-1}.}
 \tag{24}
\]
In particular, a start below $q_*$ requires
\[
 t< h\log_2\!\left(\frac32\right)
 +\log_2\!\left(
 \frac{q_*}{q_*+1-(3/2)^h}
 \right).
 \tag{25}
\]
Thus the remaining prefix problem is a low-excess inverse problem, not
an unrestricted search over exponent words.
\end{corollary}

\begin{proof}
Equation (24) is the contrapositive of (21).  Taking $N=q_*$ and
factoring
\[
 \mathcal A_h(q_*)-1
 =\left(\frac23\right)^h
 \bigl(q_*+1-(3/2)^h\bigr)
\]
gives (25).
\end{proof}

\section{The excess-zero and excess-one layers}

\begin{lemma}[first two decoder layers at a fixed suffix start]
Let $r$ be a positive odd integer and put
\[
 b:=v_3(r+1),
 \qquad s:=\frac{r+1}{3^b},
 \qquad N_1:=2^{b+1}s-1,
 \qquad b_1:=v_3(N_1).
 \tag{26}
\]
Then the following classifications are exact.

\begin{enumerate}
 \item An exact $h$-step predecessor word of excess
 $\mass(E)-h=0$ exists if and only if $0\leq h\leq b$.  It is unique:
 \[
  E=1^h,
  \qquad z+1=\left(\frac23\right)^h(r+1).
  \tag{27}
 \]

 \item Let $h>b$ and write $j:=h-b$.  An exact $h$-step predecessor
 word of excess $\mass(E)-h=1$ exists if and only if
 \[
  1\leq j\leq b_1.
  \tag{28}
 \]
 It is unique:
 \[
  \boxed{
  E=1^{j-1}2\,1^b,
  \qquad z+1=\left(\frac23\right)^jN_1.}
  \tag{29}
 \]
\end{enumerate}
\end{lemma}

\begin{proof}
Excess zero forces every exponent to equal one.  One inverse step of
exponent one replaces $x+1$ by $2(x+1)/3$.  Exactly $b$ such steps can
be reversed from $r$, proving (27).

An excess-one word has a unique exponent two.  If it occurs in position
$j$ from the start, then (10) gives
\[
 \sh(E)=2^j3^{h-j}.
 \tag{30}
\]
The integrality condition in (4) is therefore
\[
 2^{h+1-j}(r+1)\equiv3^{h-j}\pmod {3^h}.
 \tag{31}
\]
Because $h>b$, comparison of $3$-adic valuations in (31) forces
$h-j=b$.  After division by $3^b$, the remaining congruence is
\[
 2^{b+1}s\equiv1\pmod {3^j},
 \tag{32}
\]
which is equivalent to $j\leq b_1$.  This proves necessity and the
uniqueness of the proposed word.

For sufficiency, reverse the $b$ exponent-one steps nearest $r$.  The
shifted state is then $2^bs$.  Congruence (32) for $j\geq1$ makes the
next inverse exponent exactly two, and its predecessor has shifted value
$2N_1/3$.  The remaining $j-1$ inverse exponent-one steps are exact
because $3^j\mid N_1$.  Their final shifted value is
$(2/3)^jN_1$, proving (29).
\end{proof}

\begin{corollary}[fixed-seed thresholds and the next exact gate]
Let $n$ be the globally least positive odd integer with an unbounded
accelerated orbit, and suppose that the orbit of $r$ is unbounded.  Then
every word in Lemma~4 has start at least $n$.  Consequently
\[
 r+1\geq(n+1)\left(\frac32\right)^h
 \qquad(0\leq h\leq b),
 \tag{33}
\]
and
\[
 N_1\geq(n+1)\left(\frac32\right)^j
 \qquad(1\leq j\leq b_1).
 \tag{34}
\]

Apply this with $r=q_*$ in Theorem~2.  The minimum same-depth channel is
completely exhausted through its excess-zero ladder and its entire
post-ladder excess-one layer.  After the last depth covered by (27) and
(29), every prefix in that channel has excess at least two.  Hence a
proof by a longer inverse must now do at least one of the following:
\begin{enumerate}
 \item control the nonminimal suffix channels $r=q$, $r=q+2$, and
 $r\geq2^\delta(q-1)$ from (17); or
 \item control excess at least two beyond the two explicit decoder
 layers of the minimum channel.
\end{enumerate}
Any candidate start below $n$ must additionally satisfy the sharp
necessary inequality (24) with $N=n$.
\end{corollary}

\begin{proof}
Every positive predecessor of $r$ also has an unbounded forward orbit,
so its minimality gives $z\geq n$.  Substitution of (27) and (29) proves
(33) and (34).  Lemma~4 proves the exhaustion assertion, Theorem~2 gives
the three suffix channels, and Corollary~3 gives the final necessary
inequality.
\end{proof}

\end{document}
